% Part: model-theory
% Chapter: interpolation
% Section: interpolation-proof

\documentclass[../../../include/open-logic-section]{subfiles}

\begin{document}

\olfileid[fr]{mod}{int}{prf}

\olsection{Le théorème d'interpolation de Craig}

\begin{thm}[Théorème d'interpolation de Craig]
\ollabel{thm:interpol}
Pour des !!{sentence}s $!A$ et $!B$, si $\Entails !A\lif !B$, il existe
!!a{sentence}~$!C$ tel que
$\Entails !A\lif !C$ et
$\Entails !C\lif !B$.
Chaque !!{constant}, !!{function} et
!!{predicate} autre que~$\eq$ figurant
dans~$!C$ figure à la fois dans~$!A$
et dans~$!B$. On appelle~$!C$ un
\emph{interpolant} de~$!A$ et~$!B$.\footnote{Précision éditoriale : nous explicitons que $!A$ et $!B$ sont des énoncés fermés, comme l'exige la preuve de ce chapitre. La version générale pour des formules ouvertes fournit un interpolant qui peut lui aussi avoir des variables libres.}
\end{thm}

\begin{proof}
Soient $\Lang{L}_1$ le langage de~$!A$,
$\Lang{L}_2$ celui de~$!B$, et
$\Lang{L}_0=\Lang{L}_1\cap\Lang{L}_2$.
Pour $i\in\{0,1,2\}$, formons
$\Lang{L}'_i$ en ajoutant à
$\Lang{L}_i$ les mêmes constantes
nouvelles $\Obj c_0,\Obj c_1,\Obj c_2,\dots$.

Si~$!A$ est insatisfaisable, la
contradiction $\lforall[x][\lnot\eq[x][x]]$
est un interpolant. Si~$\lnot !B$ est
insatisfaisable, la tautologie
$\lexists[x][\eq[x][x]]$ en est un.
Supposons donc $!A$ et~$\lnot !B$
satisfaisables. Pour démontrer la
contraposée du théorème, supposons
qu'il n'existe aucun interpolant.
Autrement dit, $\{!A\}$ et
$\{\lnot !B\}$ sont inséparables
dans~$\Lang{L}_0$.

Nous allons étendre cette paire à
une paire maximalement inséparable
$(\Gamma^*,\Delta^*)$.
Énumérons les !!{sentence}s de
$\Lang{L}'_1$ sous la forme
$!A_0,!A_1,!A_2,\dots$, et ceux de
$\Lang{L}'_2$ sous la forme
$!B_0,!B_1,!B_2,\dots$.
Choisissons des énumérations dans
lesquelles chaque !!{sentence}
revient une infinité de fois.\footnote{Compléments éditoriaux : les énumérations portent sur les langages étendus, alors que la source écrit les langages initiaux. Les répétitions garantissent qu'un énoncé existentiel ajouté après son premier passage reçoit ensuite un témoin.}
Posons $\Gamma_0=\{!A\}$ et
$\Delta_0=\{\lnot !B\}$, puis
construisons par récurrence des paires
croissantes $(\Gamma_n,\Delta_n)$.
À l'étape~$n$ :
\begin{enumerate}
\item Si $\Gamma_n\cup\{!A_n\}$ et
  $\Delta_n$ sont inséparables dans
  $\Lang{L}'_0$, ajoutons $!A_n$ à
  $\Gamma_{n+1}$. Si $!A_n$ a la forme
  $\lexists[x][!S]$, ajoutons aussi
  $\Subst{!S}{c}{x}$, où~$c$ est
  une constante de la famille commune $\Obj c_j$, absente de
  $\Gamma_n$, $\Delta_n$, $!A_n$
  et~$!B_n$. Si la paire $\Gamma_n\cup\{!A_n\}$, $\Delta_n$ est séparable,
  posons $\Gamma_{n+1}=\Gamma_n$.
\item Si $\Gamma_{n+1}$ et
  $\Delta_n\cup\{!B_n\}$ sont
  inséparables dans $\Lang{L}'_0$,
  ajoutons $!B_n$ à $\Delta_{n+1}$.
  Si $!B_n=\lexists[x][!S]$, ajoutons
  également $\Subst{!S}{d}{x}$ pour
  une constante~$d$ de la famille commune $\Obj c_j$, absente
  de $\Gamma_{n+1}$, $\Delta_n$,
  $!A_n$ et~$!B_n$. Si la paire $\Gamma_{n+1}$, $\Delta_n\cup\{!B_n\}$ est séparable, posons
  $\Delta_{n+1}=\Delta_n$.
\end{enumerate}
Enfin, posons
\begin{align*}
  \Gamma^*&=\bigcup_{n\geq0}\Gamma_n,
  &
  \Delta^*&=\bigcup_{n\geq0}\Delta_n.
\end{align*}

Vérifions simultanément par récurrence
les deux affirmations suivantes :
\begin{enumerate}
\item\ollabel{part-a}
  $\Gamma_n$ et~$\Delta_n$ sont
  inséparables pour tout~$n$;
\item\ollabel{part-b}
  $\Gamma_{n+1}$ et~$\Delta_n$
  sont inséparables pour tout~$n$.
\end{enumerate}
Le cas initial résulte du
\olref[sep]{lem:sep1}. À chaque
demi-étape, si l'énoncé considéré
est rejeté, la paire ne change pas.
S'il est ajouté sans témoin, elle
reste inséparable par le critère
de l'étape. S'il est existentiel,
le \olref[sep]{lem:sep2} garantit
qu'elle reste inséparable après
l'adjonction du témoin frais.
Cela prouve \olref{part-a} et
\olref{part-b} à toutes les étapes.
Si un !!{sentence} séparait
$\Gamma^*$ et~$\Delta^*$, la
compacité donnerait des sous-ensembles
finis séparés, contenus dans
$\Gamma_n$ et~$\Delta_n$ pour un
même~$n$. C'est impossible.
Chacun des deux ensembles limites
est donc cohérent : s'il était
incohérent, une contradiction ou
une tautologie du langage commun
séparerait la paire.

Montrons que $\Gamma^*$ est
maximalement cohérent dans
$\Lang{L}'_1$. Supposons qu'un
!!{sentence}~$!E$ et sa négation
soient tous deux absents de
$\Gamma^*$. Prenons une étape~$n$
où $!A_n=!E$. Comme $!E$ n'a pas
été ajouté, il existe un
!!{sentence}~$!C$ de~$\Lang{L}'_0$
tel que
\begin{align*}
  \Gamma_n\cup\{!E\}&\Entails !C,
  &
  \Delta_n&\Entails\lnot !C.
\end{align*}
Ainsi $\Gamma^*\Entails !E\lif !C$
et $\Delta^*\Entails\lnot !C$.
À une étape où $\lnot !E$ est
considéré et rejeté, on obtient
de même un $!C'$ tel que
$\Gamma^*\Entails\lnot !E\lif !C'$
et $\Delta^*\Entails\lnot !C'$.
Alors $\Gamma^*\Entails !C\lor !C'$,
tandis que
$\Delta^*\Entails\lnot(!C\lor !C')$,
contradiction. Le même raisonnement,
en échangeant les deux côtés,
montre que $\Delta^*$ est
maximalement cohérent dans
$\Lang{L}'_2$.

Leur intersection est maximalement
cohérente dans~$\Lang{L}'_0$.
Elle est cohérente, puisqu'elle
est contenue dans chacun des deux
ensembles cohérents. Pour tout
!!{sentence} commun~$!S$, chacun
des deux ensembles limites contient
exactement l'un de $!S$ et~$\lnot !S$.
Ils ne peuvent faire des choix
opposés : $!S$ ou~$\lnot !S$
les séparerait. Ils font donc
le même choix, qui appartient
à leur intersection.

Les répétitions dans nos
énumérations assurent aussi la
propriété des témoins. Si
$\lexists[x][!S]$ entre dans
$\Gamma^*$, il est de nouveau
considéré après cette entrée;
un témoin $\Subst{!S}{c}{x}$
est alors ajouté. Il en va de
même pour~$\Delta^*$. Prenons
les modèles canoniques de termes
$\Struct{M}'_1$ et~$\Struct{M}'_2$
de ces deux théories complètes et
pourvues de témoins, comme dans
la preuve du théorème de complétude
(\olref[fol][com][cth]{thm:completeness}).
Tout élément de chacun de ces
modèles est nommé par une constante de la famille commune $\Obj c_j$ :
pour tout terme clos~$t$, l'!!{sentence} $\lexists[x][\eq[x][t]]$
est valide et obtient finalement
un témoin constant égal à~$t$.

Réduisons $\Struct{M}'_1$ et
$\Struct{M}'_2$ au langage commun
$\Lang{L}'_0$, et notons les
structures obtenues $\Struct{M_1}$
et~$\Struct{M_2}$. L'application
\[
  h\colon\Domain{M_1}\longrightarrow
  \Domain{M_2},\qquad
  h(\Assign{c}{M'_1})=\Assign{c}{M'_2}
\]
pour $c$ dans la famille commune $\Obj c_j$, est bien définie : une égalité
entre constantes est vraie dans
un modèle si et seulement si elle
appartient à sa théorie complète,
et les deux théories s'accordent
sur tous les !!{sentence}s de
$\Lang{L}'_0$. Le même argument
donne l'injectivité et, puisque
tout élément est nommé, la
surjectivité. Les !!{sentence}s
atomiques communes portant sur
des constantes montrent que~$h$
préserve les !!{predicate}s et
les !!{function}s. C'est donc
un isomorphisme des structures
réduites.

Construisons une structure
$\Struct{M}$ du langage
$\Lang{L}'_1\cup\Lang{L}'_2$
sur le domaine $\Domain{M_2}$.
Pour les symboles de
$\Lang{L}'_2$, conservons leur
interprétation dans
$\Struct{M}'_2$. Pour un
!!{predicate}~$P$ de
$\Lang{L}'_1\setminus\Lang{L}'_2$,
transportons l'interprétation
de~$\Struct{M}'_1$ par~$h$ :
\[
  \tuple{h(a_1),\dots,h(a_k)}
  \in\Assign{P}{M}
  \quad\Longleftrightarrow\quad
  \tuple{a_1,\dots,a_k}
  \in\Assign{P}{M'_1}.
\]
Pour un !!{function}~$f$
propre à~$\Lang{L}'_1$, posons
\[
\Assign{f}{M}(h(a_1),\dots,h(a_k))
=h\bigl(\Assign{f}{M'_1}(a_1,\dots,a_k)\bigr).
\]
La même règle couvre les
constantes. Sur les symboles
communs, les deux prescriptions
coïncident puisque~$h$ est un
isomorphisme.\footnote{Corrections éditoriales : les symboles propres à $\Lang{L}'_1$ sont transportés depuis $\Struct{M}'_1$, et non depuis $\Struct{M}'_2$ comme l'écrit par erreur la source pour les prédicats. La structure réunie est construite dans le langage étendu, puisqu'elle doit satisfaire les théories contenant les constantes nouvelles.} Par récurrence
sur les !!{formula}s,
$\Struct{M}$ satisfait les mêmes
!!{sentence}s de $\Lang{L}'_1$
que~$\Struct{M}'_1$ et les mêmes
!!{sentence}s de $\Lang{L}'_2$
que~$\Struct{M}'_2$. Ainsi
$\Struct{M}\Entails\Gamma^*\cup\Delta^*$,
en particulier
$\Struct{M}\Entails !A$ et
$\Struct{M}\Entails\lnot !B$.
Cela contredit $\Entails !A\lif !B$.
\end{proof}

\end{document}
